% Proof of the canonical-collar repair of the source-legal arbitrary-field
% virtual-annulus comparison.
% Date: 28 August 2026.
% Standalone production-verification candidate; not integrated into the book.
% The rejected c5a643/292bb and 0315d3/c903cf artifacts are preserved
% unchanged in separate files.

\begin{proof}[Proof of
\cref{lem:p21p3-gf-legal-source-localization}]
Use the ordinary local moment description in
\cite[Proposition~15.1.1(b)]{CLD25}.  After shrinking \(W\) about \(y\),
there are an analytic moment \(F_{\mathrm{form}}\) and a smooth superform
\(\eta\) such that \(\alpha=F_{\mathrm{form}}^*\eta\).  Shrink once more so that
the prescribed finite block \(G\) is defined.  When \(\chi\) is present,
\cite[Corollary~15.1.2]{CLD25} gives, after one further shrinking, an analytic
moment \(H\) and a smooth function \(\Phi\) such that
\(\chi=\Phi\circ H_{\mathrm{trop}}\).  Relabel this common neighborhood as
\(W\).

Apply \cite[Lemma~12.3.6]{CLD25} to the analytic space \(W\) at \(y\).  It
gives a compact analytic neighborhood \(V\) of \(y\) in \(W\) and a moment
\(F_{\mathrm{bd}}\) such that
\[
 F_{\mathrm{bd,trop}}(y)\notin
 F_{\mathrm{bd,trop}}(\partial(V/W)).
\]
Concatenate \(F_{\mathrm{bd}}\), \(F_{\mathrm{form}}\), \(G\), and \(H\).
Projection to the first block proves
\eqref{eq:p21p3-gf-source-local-boundary}; projection to the form and cutoff
blocks proves \eqref{eq:p21p3-gf-source-local-form} and
\eqref{eq:p21p3-gf-source-cutoff-factorization}.  Equality of two full
tropical values implies equality of their \(H\)-projections, so applying
\(\Phi\) proves \eqref{eq:p21p3-gf-source-cutoff-fibre-constancy}.

The remainder is the localization used in the proof of
\cite[Theorem~16.8.5]{CLD25}.  The set \(\Omega\) in
\eqref{eq:p21p3-gf-source-local-omega} is open because the relative boundary
is compact.  The set \(U\) misses \(\partial(V/W)\), hence is contained in the
relative interior of \(V\) and is open in \(W\).  It contains \(y\) by
\eqref{eq:p21p3-gf-source-local-boundary}.  Moreover
\(F_{\mathrm{trop}}|_U\) is obtained from the proper map
\(F_{\mathrm{trop}}:V\to\mathbf R^N\) by base change along
\(\Omega\hookrightarrow\mathbf R^N\), and is therefore proper.  Thus the
inverse image of the compact support of \(\lambda\) is compact in \(U\).
The displayed pullback form consequently extends by zero across
\(W\setminus U\), and the equality with \(\alpha\) where \(\lambda=1\)
follows from \eqref{eq:p21p3-gf-source-local-form}.

For the last assertion, apply the construction at each point of the compact
test support, starting inside the prescribed member of the source cover.
The resulting open sets \(U\) cover that support, so finitely many suffice.
Smooth cutoffs subordinate to this finite cover, with sum one near the
compact support, exist by \cite[Corollary~15.2.9]{CLD25}.  Current identities
are local and therefore glue after this finite decomposition.  If one of
those cutoffs is retained in a
single local presentation, the preceding application of
\cite[Corollary~15.1.2]{CLD25} supplies \(H\); at no point was the cutoff
itself used as an analytic coordinate.
\end{proof}

\begin{proof}[Proof of
\cref{prop:p21p3-gf-virtual-annulus-bridge}]
The residual curve and the branch--valuation correspondence over an
arbitrary complete field are
\cite[\S\S3.3.5.2, 4.2.6, 4.2.9, and 4.2.11.1]
{Ducros24CurveStructure}.  Since \(x\)
is interior, every closed point of \(C_x\) occurs as a branch.  For the branch
\(b\) corresponding to \(p_b\), the refined value group is the lexicographic
extension of \(|\mathcal H(x)^{\times}|\) by the discrete valuation
\(\operatorname{ord}_{p_b}\).  If \(s_b\) is a uniformizer of the discrete
valuation ring \(\mathcal O_{C_x,p_b}\), then
\cite[\S4.2.11.1]{Ducros24CurveStructure} says precisely that the class of
its refined value generates
\[
 |\mathcal H(x)^{\times}|_b/|\mathcal H(x)^{\times}|\simeq\mathbf Z.
\]
Since \(Y\) is normal and \(x\) is nonrigid,
\cite[\S4.2.18]{Ducros24CurveStructure} identifies
\(\mathcal O_{Y,x}\) with \(\kappa(x)\), and
\cite[\S4.2.6]{Ducros24CurveStructure} says that \(\kappa(x)\) is dense in
\(\mathcal H(x)\).  Hence \(s_b\), viewed as a homogeneous residue element
of degree one, has a norm-one lift \(f_b\in\mathcal O_{Y,x}\).  This is also
the finite residue-lifting step made explicit in the proof of
\cite[Lemma~6.14]{GJR26}.  Its value in \(\mathcal H(x)\) is nonzero, so
after shrinking it is an analytic unit.  This proves that
\eqref{eq:p21p3-gf-branch-order-definition} is meaningful.  Replacing a
uniformizer by another one multiplies it by a unit in
\(\mathcal O_{C_x,p_b}\), whose class in the displayed quotient is trivial.
This proves independence and also gives
\eqref{eq:p21p3-gf-branch-order-residue}.

We next prove the exact ray assertion.  By
\cite[Proposition~4.3.9(2) and \S\S4.3.10.2--4.3.10.3]
{Ducros24CurveStructure}, \(b\) has a cofinal family of virtual-annular
sections.  Choose one and call it \(Z\).  Its analytic skeleton and canonical
retraction are compatible with every complete scalar extension by
\cite[\S\S3.6.35.5--3.6.35.10]{Ducros24CurveStructure}.  After a finite
separable extension and then passage to a completed algebraic closure, a
component above \(Z\) is an ordinary annulus whose skeleton maps
homeomorphically to \(S^{\mathrm{an}}(Z)\).

Let \(\beta\) be the branch of this split annulus above \(b\), and choose an
annulus coordinate \(\theta\) so that \(-\log|\theta|\) increases away from
the endpoint over \(x\).
The annular calculation in
\cite[\S\S3.6.22--3.6.24 and \S4.3.6.4]
{Ducros24CurveStructure} shows that, after shortening the annulus once for
the finite family \(g_1,\ldots,g_m\),
\begin{equation}
\label{eq:p21p3-gf-split-annulus-monomials}
                  |g_j|=c_j|\theta|^{n_j}
                  \qquad(1\leq j\leq m)
\end{equation}
for constants \(c_j>0\) and integers \(n_j\).  Naturality of the refined
branch valuation is \cite[\S4.2.9]{Ducros24CurveStructure}.  Consequently,
restriction from \(b\) to \(\beta\) sends the generator defined by \(g_m\) to
a positive integral multiple of the primitive annular generator.  Hence
there is one integer \(e_b>0\), independent of \(j\), such that
\begin{equation}
\label{eq:p21p3-gf-common-expansion-index}
                       n_j=e_b\partial_b(g_j)
                       \qquad(1\leq j\leq m).
\end{equation}
Indeed, this is just the homomorphism between the two rank-one quotient
value groups; \(n_m=e_b\) because \(\partial_b(g_m)=1\).  The same integer
is the tangent-lattice expansion index of this split annular presentation
by \cite[\S3.6.24.2]{Ducros24CurveStructure}.  In the finite-morphism
setting of \cite[Theorem~4.3.13(i)]{Ducros24CurveStructure}, that annular
index is identified with the corresponding residual ramification index.

Put
\[
 t=-\log|g_m|+\log|g_m(x)|.
\]
The orientation and
\cite[\S4.2.11.1]{Ducros24CurveStructure} make \(t\) strictly positive
away from \(x\) on a sufficiently short section.  Equations
\eqref{eq:p21p3-gf-split-annulus-monomials} and
\eqref{eq:p21p3-gf-common-expansion-index}, followed by passage to the
limit at \(x\), give on the skeleton
\begin{equation}
\label{eq:p21p3-gf-coordinate-ray-formula}
 -\log|g_j|+\log|g_j(x)|
       =\partial_b(g_j)t
       \qquad(1\leq j\leq m).
\end{equation}
Every invertible analytic function on a virtual annulus has the same norm
as at its canonical retraction by
\cite[\S3.6.35.8]{Ducros24CurveStructure}.  Thus
\eqref{eq:p21p3-gf-coordinate-ray-formula} holds on the whole shortened
section \(Z_b\), and \(t\) identifies its skeleton with an interval
\((0,\varepsilon_b)\).  This proves
\eqref{eq:p21p3-gf-exact-branch-ray}.  Since the last coordinate of \(v_b\)
is one, \(v_b\) is primitive in the saturated ambient lattice.  The
one-dimensional skeleton attached to the moment, restricted to any compact
subannular band, has the unique integral affine structure for which the
moment maps its faces integral-affinely onto their images by
\cite[\S8.13]{GJR26}; hence this is also the primitive tangent of the
corresponding source one-cell.  This proves
\cref{itm:p21p3-gf-bridge:ray}.

We prove the finite separation assertion.  Let \(p_i=p_{b_i}\).  Weak
approximation on the normal projective curve \(C_x\) gives rational functions
\(s_1,\ldots,s_s\) with
\begin{equation}
\label{eq:p21p3-gf-kronecker-orders}
                \operatorname{ord}_{p_i}(s_j)=\delta_{ij}.
\end{equation}
For completeness, choose a closed point \(q\) away from the \(p_i\).  For
\(n\gg0\), Serre vanishing gives
\[
 H^1\!\left(C_x,
   \mathcal O_{C_x}\!\left(nq-2\sum_i p_i\right)\right)=0.
\]
The divisor exact sequence makes the simultaneous value-and-first-jet map
from \(H^0(C_x,\mathcal O_{C_x}(nq))\) onto
\(\bigoplus_i\mathcal O_{C_x,p_i}/\mathfrak m_{p_i}^2\) surjective.
Prescribing a nonzero first jet and value zero at \(p_j\), and value one at
the other \(p_i\), yields \(s_j\) satisfying
\eqref{eq:p21p3-gf-kronecker-orders}.  This argument is valid for closed
points of arbitrary residue degree.

Lift the \(s_j\) to norm-one analytic units by the density argument above
(equivalently, the lifting step in the proof of
\cite[Lemma~6.14]{GJR26}), and first append those lifts to \(F_0\).  By
\cite[Proposition~4.2.4(ii)]{Ducros24CurveStructure}, every
analytic unit has nonzero branch order in only finitely many branches.
The preliminary outgoing vector in branch \(b_i\) therefore has, in the new
\(s_1,\ldots,s_s\) block, the standard basis vector \(e_i\); in particular it
is primitive, and the selected vectors are pairwise nonparallel.

There can still be an unselected branch whose preliminary vector is a
positive multiple of the vector for some \(b_i\).  There are only finitely
many such branches: any one of them has positive \(s_i\)-order and hence
belongs to the finite exceptional set supplied by
\cite[Proposition~4.2.4(ii)]{Ducros24CurveStructure}.  Call their union
\(Q\).  For every \(c\in Q\), weak approximation as above gives a rational
function \(h_c\) satisfying
\begin{equation}
\label{eq:p21p3-gf-collision-killers}
 \operatorname{ord}_{p_{b_i}}(h_c)=0\quad(1\leq i\leq s),
 \qquad
 \operatorname{ord}_{p_c}(h_c)=1.
\end{equation}
Lift all \(h_c\) and append them as well.  The new coordinates vanish in
branch order on every selected branch, whereas the \(h_c\)-coordinate has
order one on \(c\).  Thus each old collision is destroyed.  A newly active
branch cannot create a new collision: projection to the preliminary
coordinates would already have made it an element of \(Q\).  This proves
\cref{itm:p21p3-gf-bridge:separation}.

We finally prove \cref{itm:p21p3-gf-bridge:localization}.  First use
\cref{lem:p21p3-gf-legal-source-localization} inside the prescribed source
\(W\), with no branch block yet, and retain the resulting moment germ
\(F_0\).  Apply
\cref{itm:p21p3-gf-bridge:separation} to \(F_0\).  On a common neighborhood
of \(x\), let \(F_{\mathrm{br}}\) be the resulting separating block and put
\(G_0=(F_0,F_{\mathrm{br}})\).  Let \(O\subset W\) be a common open
neighborhood of \(x\) on which every coordinate of \(G_0\) is an analytic
unit.  We choose the source star before adding any further moment coordinates.
By \cite[Theorem~4.5.4]{Ducros24CurveStructure}, after shrinking inside
\(O\), there is a connected open elementary neighborhood \(V_*\) of \(x\)
whose closure is a compact tree and for which
\begin{equation}
 V_*\setminus\{x\}
   =\coprod_{\lambda\in\Lambda}B_\lambda
      \amalg\coprod_{\nu=1}^{r}A_\nu .
 \label{eq:p21p3-gf-elementary-star-decomposition}
\end{equation}
The \(B_\lambda\) are virtual discs, whereas
\(A_1,\ldots,A_r\) are the finitely many gently virtual-annular components.
Every component in \eqref{eq:p21p3-gf-elementary-star-decomposition}
represents one branch at \(x\), and the compact-tree condition prevents an
annular component from returning to \(x\) at its other end.

An invertible analytic function has constant norm on a virtual disc by
\cite[\S3.6.32.3]{Ducros24CurveStructure}.  Continuity at \(x\) therefore
gives
\begin{equation}
 G_{0,\mathrm{trop}}|_{B_\lambda}
       =G_{0,\mathrm{trop}}(x)
       \qquad(\lambda\in\Lambda).
 \label{eq:p21p3-gf-disc-components-at-vertex}
\end{equation}
This retains all, possibly infinitely many, virtual-disc components whole.

Only the finitely many annular components need to be shortened.  Let
\(b_\nu\) be the branch represented by \(A_\nu\), and put
\(w_\nu=(\partial_{b_\nu}(g))_g\), where \(g\) runs through the coordinates
of \(G_0\).  For this calculation only, choose a norm-one analytic unit
\(u_\nu\in\mathcal O_{Y,x}^{\times}\) with
\(\partial_{b_\nu}(u_\nu)=1\); it is not appended to the moment.  Since there
are only finitely many \(A_\nu\), all the \(u_\nu\) are units on a common
sufficiently short collar.  After a finite separable extension splitting
\(A_\nu\), choose an annulus coordinate \(\theta_\nu\) oriented away from the
end over \(x\).  Apply \eqref{eq:p21p3-gf-split-annulus-monomials}
simultaneously to the coordinates of \(G_0\) and to \(u_\nu\).  Naturality of
the refined branch valuation, \cite[\S4.2.9]{Ducros24CurveStructure}, and the
rank-one quotient argument of
\eqref{eq:p21p3-gf-common-expansion-index} give one integer \(e_\nu>0\) such
that the split-annulus exponents are
\[
 n_{\nu,g}=e_\nu\partial_{b_\nu}(g),
 \qquad n_{\nu,u}=e_\nu.
\]
Consequently, after choosing a skeleton point \(q_\nu\) sufficiently close to
the end accumulating at \(x\),
\begin{equation}
 G_{0,\mathrm{trop}}(z)
   =G_{0,\mathrm{trop}}(x)+t_\nu(z)w_\nu,
 \qquad 0<t_\nu(z)<\delta_\nu,
 \label{eq:p21p3-gf-finite-annulus-ray-formula}
\end{equation}
on the retained collar, where
\[
 t_\nu(z)=-\log|u_\nu(z)|+\log|u_\nu(x)|.
\]
The orientation makes \(t_\nu\) positive and strictly increasing along the
retained skeleton.  By \cite[\S3.6.35.8]{Ducros24CurveStructure} it factors
through the canonical retraction.  If \(w_\nu=0\), the displayed formula says
that the whole retained collar maps to the vertex.
Sections~3.6.35.9--3.6.35.10 of \cite{Ducros24CurveStructure} identify the
inverse image \(A_\nu^0\) of the retained skeleton interval with an analytic
virtual subannulus, compatibly with scalar extension.

Let \(W_1\) be the component containing \(x\) of
\(V_*\setminus\{q_1,\ldots,q_r\}\).  Analytic curves are locally connected
graphs by \cite[Theorem~3.5.1]{Ducros24CurveStructure}; hence \(W_1\) is an
open connected analytic neighborhood of \(x\).  The tree description and
the choice of the \(q_\nu\) give
\begin{equation}
 W_1\setminus\{x\}
   =\coprod_{\lambda\in\Lambda}B_\lambda
      \amalg\coprod_{\nu=1}^{r}A_\nu^0.
 \label{eq:p21p3-gf-trimmed-star-decomposition}
\end{equation}
Thus no independent truncation has been made in infinitely many directions.

Equation \eqref{eq:p21p3-gf-disc-components-at-vertex} and the order-one
separator coordinate show that a selected branch cannot be represented by a
virtual-disc component.  Hence it is represented by a unique
\(A_{\nu(i)}^0\).  Put \(w_i=w_{\nu(i)}\) and
\(\delta_i=\delta_{\nu(i)}\).  The residue-unit block contains a coordinate
of order one on \(b_i\), so \(w_i\ne0\), and
\cref{itm:p21p3-gf-bridge:separation} says that
\(\mathbf R_{>0}w_i\) is different from the ray of every other active
component.  After moving the finitely many \(q_\nu\) closer to \(x\), if
necessary, equations \eqref{eq:p21p3-gf-disc-components-at-vertex} and
\eqref{eq:p21p3-gf-finite-annulus-ray-formula} therefore give the exact
preliminary locator
\begin{equation}
 G_{0,\mathrm{trop}}^{-1}\!\left(
 G_{0,\mathrm{trop}}(x)+(0,\delta_i)w_i\right)\cap W_1
 = A_{\nu(i)}^0.
\end{equation}
Indeed, every virtual disc and every inactive annulus maps to the vertex,
every other active annulus maps into a distinct open ray, and on the selected
annulus the order-one coordinate makes \(t_\nu\) a genuine radial parameter.
This excludes remote source collisions without an infinite simultaneous
choice of branch sections.

Reapply \cref{lem:p21p3-gf-legal-source-localization} inside \(W_1\), with
\(G=G_0\), and call its compact domain and full moment \(V\) and \(F\).  The
lemma gives boundary avoidance for this full moment.  Hence the open set
\(U\) in \eqref{eq:p21p3-gf-source-local-omega} contains \(x\), lies in
\(W_1\), and \(F_{\mathrm{trop}}|_U\) is proper.  Let
\(P:\mathbf R^N\to\mathbf R^{N_0}\) be projection to the \(G_0\)-block, so
that \(P\circ F_{\mathrm{trop}}=G_{0,\mathrm{trop}}\).

For a selected branch \(b_i\), the full moment \(F\) still contains the
order-one separator coordinate.  On \(A_{\nu(i)}^0\), the formula above for
that coordinate identifies its normalized negative logarithm with the radial
parameter \(t_i=t_{\nu(i)}\).  Since \(U\) is an open neighborhood of \(x\),
\cite[\S4.3.10.3]{Ducros24CurveStructure} and
\cite[\S\S3.6.35.9--3.6.35.10]{Ducros24CurveStructure} allow us to choose
\(0<\varepsilon_i<\delta_i\) so that
\begin{equation}
 Z_{b_i}:=t_i^{-1}((0,\varepsilon_i))
       \subset A_{\nu(i)}^0\cap U
 \label{eq:p21p3-gf-canonical-full-source-collar}
\end{equation}
is the full canonical virtual subannulus over that skeleton interval.  After
decreasing \(\varepsilon_i\) once more, the calculation in
\cref{itm:p21p3-gf-bridge:ray}, applied simultaneously to all coordinates of
\(F\), gives
\[
 F_{\mathrm{trop}}(z)=F_{\mathrm{trop}}(x)+t_i(z)v_i
 \qquad(z\in Z_{b_i}),
\]
where \(v_i\) is the full primitive branch-order vector.  This proves
\eqref{eq:p21p3-gf-source-separated-ray-image}; it also gives
\(P(v_i)=w_i\).  If \(z\in U\) and
\(F_{\mathrm{trop}}(z)=F_{\mathrm{trop}}(x)+t v_i\) for
\(0<t<\varepsilon_i\), projection and
\eqref{eq:p21p3-gf-preliminary-exact-source-locator} put \(z\) in
\(A_{\nu(i)}^0\).  The order-one separator coordinate then gives
\(t_i(z)=t\).  By the definition
\eqref{eq:p21p3-gf-canonical-full-source-collar}, \(z\in Z_{b_i}\).  The
reverse inclusion follows from the displayed ray formula, and hence
\eqref{eq:p21p3-gf-source-separated-ray-preimage} holds.

Consequently the Euclidean current
\[
 \langle T_{\Omega},\beta\rangle
       =\langle T|_U,F^*\beta\rangle
\]
used in the proof of \cite[Theorem~16.8.5]{CLD25}, when tested on a
sufficiently short relative-open subsegment of the selected ray, receives
source contribution only from \(Z_{b_i}\).  If a norm-one lift of
\(g\in\operatorname{Rat}(C_x)^\times\) is appended, projection to the
previous coordinates preserves the exact separation, while
\eqref{eq:p21p3-gf-branch-order-residue} identifies its scalar slope with
\(\operatorname{ord}_{p_{b_i}}(g)\).  Multiplication by the edge flux \(a_i\)
gives the asserted balancing term.  No smooth plateau was used as a moment
coordinate, and no scalar extension of \(T\) occurred.

Equation \eqref{eq:p21p3-gf-common-expansion-index} also proves the claim in
\cref{rem:p21p3-gf-virtual-annulus-index}.  Under pushforward, the rank-one
matrix current has one factor \(e_b\): its tangent tensor contributes
\(e_b^2\) and one-dimensional change of measure contributes \(e_b^{-1}\).
The canonical calibration has the same expansion factor by
\cite[\S14.3.2, Lemma~14.3.4, and Proposition~14.4.2]{CLD25}.  No arbitrary
current was extended to a splitting field at any point of the argument.
\end{proof}
